% GENERATED FILE. Edit canonical lessons, then run npm run generate:books.
% canonical-source-hash: fnv1a64:696bfa0b28aabf57
% canonical-lessons: JA-C107-doro, JA-C107-hokori, JA-C107-awa, JA-C107-tsuyu, JA-C107-saka

\chapter{Mud, Dust, Foam, Dew, Slope}
\label{ch:ja-mud-dust-foam-dew-slope}
\hlchaptermodality{107}

\begin{chapteropening}
\textbf{By the end of this chapter:} \emph{I can say mud, dust, foam, dew and a slope.}

Complete the last lesson of chapter 107: I can say mud, dust, foam, dew and a slope.
\end{chapteropening}

\section[\texorpdfstring{\ja{どろ} (doro)}{どろ (doro)}]{\ja{どろ} (doro) — mud}

\label{lesson:JA-C107-doro}



\begin{quote}
Before the new one: say the Japanese for a rock, then the Japanese for a pebble.
\end{quote}

\subsection*{You'll want to know: \ja{どろ}}
\textbf{\ja{どろ}} — \emph{doro} — \textquotedblleft{}mud\textquotedblright{}.

\begin{grammarlens}[title={What you've built}]
One more word for this chapter.
\end{grammarlens}

\subsection*{Your turn}
\begin{itemize}
\raggedright
  \item \emph{Say it:} \emph{doro}
  \item \emph{Say it:} \emph{doro}, once more
  \item \emph{Say it:} say \emph{doro}
  \item \emph{Recall:} say the Japanese for a rock, then the Japanese for a pebble, then say \emph{doro} again
\end{itemize}

\begin{tcolorbox}[breakable,colback=teal!4,colframe=teal!35!black,title={Before you move on}]
What does \ja{どろ} mean? (\textquotedblleft{}Mud\textquotedblright{}.) Say it once more.
\end{tcolorbox}
\section[\texorpdfstring{\ja{ほこり} (hokori)}{ほこり (hokori)}]{\ja{ほこり} (hokori) — dust}

\label{lesson:JA-C107-hokori}



\begin{quote}
Before the new one: say the Japanese for a pebble, then the Japanese for mud.
\end{quote}

\subsection*{You'll want to know: \ja{ほこり}}
\textbf{\ja{ほこり}} — \emph{hokori} — \textquotedblleft{}dust\textquotedblright{}.

\begin{grammarlens}[title={What you've built}]
One more word for this chapter.
\end{grammarlens}

\subsection*{Your turn}
\begin{itemize}
\raggedright
  \item \emph{Say it:} \emph{hokori}
  \item \emph{Say it:} \emph{hokori}, once more
  \item \emph{Say it:} say \emph{hokori}
  \item \emph{Recall:} say the Japanese for a pebble, then the Japanese for mud, then say \emph{hokori} again
\end{itemize}

\begin{tcolorbox}[breakable,colback=teal!4,colframe=teal!35!black,title={Before you move on}]
What does \ja{ほこり} mean? (\textquotedblleft{}Dust\textquotedblright{}.) Say it once more.
\end{tcolorbox}
\section[\texorpdfstring{\ja{あわ} (awa)}{あわ (awa)}]{\ja{あわ} (awa) — foam, bubbles}

\label{lesson:JA-C107-awa}



\begin{quote}
Before the new one: say the Japanese for mud, then the Japanese for dust.

Read each word whole: \textbf{\ja{みみ}}, \textbf{\ja{くち}}, \textbf{\ja{かお}}.
\end{quote}

\subsection*{You'll want to know: \ja{あわ}}
\textbf{\ja{あわ}} — \emph{awa} — \textquotedblleft{}foam, bubbles\textquotedblright{}.

\begin{grammarlens}[title={What you've built}]
One more word for this chapter.
\end{grammarlens}

\subsection*{Your turn}
\begin{itemize}
\raggedright
  \item \emph{Say it:} \emph{awa}
  \item \emph{Say it:} \emph{awa}, once more
  \item \emph{Say it:} say \emph{awa}
  \item \emph{Recall:} say the Japanese for mud, then the Japanese for dust, then say \emph{awa} again
\end{itemize}

\begin{tcolorbox}[breakable,colback=teal!4,colframe=teal!35!black,title={Before you move on}]
What does \ja{あわ} mean? (\textquotedblleft{}Foam, bubbles\textquotedblright{}.) Say it once more.
\end{tcolorbox}
\section[\texorpdfstring{\ja{つゆ} (tsuyu)}{つゆ (tsuyu)}]{\ja{つゆ} (tsuyu) — dew; the rainy season}

\label{lesson:JA-C107-tsuyu}



\begin{quote}
Before the new one: say the Japanese for dust, then the Japanese for foam.
\end{quote}

\subsection*{You'll want to know: \ja{つゆ}}
\textbf{\ja{つゆ}} — \emph{tsuyu} — \textquotedblleft{}dew; the rainy season\textquotedblright{}.

\begin{grammarlens}[title={What you've built}]
One more word for this chapter.
\end{grammarlens}

\subsection*{Your turn}
\begin{itemize}
\raggedright
  \item \emph{Say it:} \emph{tsuyu}
  \item \emph{Say it:} \emph{tsuyu}, once more
  \item \emph{Say it:} say \emph{tsuyu}
  \item \emph{Recall:} say the Japanese for dust, then the Japanese for foam, then say \emph{tsuyu} again
\end{itemize}

\begin{tcolorbox}[breakable,colback=teal!4,colframe=teal!35!black,title={Before you move on}]
What does \ja{つゆ} mean? (\textquotedblleft{}Dew; the rainy season\textquotedblright{}.) Say it once more.
\end{tcolorbox}
\section[\texorpdfstring{\ja{さか} (saka)}{さか (saka)}]{\ja{さか} (saka) — a slope}

\label{lesson:JA-C107-saka}



\begin{quote}
Before the new one: say the Japanese for foam, then the Japanese for dew; the rainy season.
\end{quote}

\subsection*{You'll want to know: \ja{さか}}
\textbf{\ja{さか}} — \emph{saka} — \textquotedblleft{}a slope\textquotedblright{}.

\begin{grammarlens}[title={What you've built}]
One more word for this chapter.
\end{grammarlens}

\subsection*{Your turn}
\begin{itemize}
\raggedright
  \item \emph{Say it:} \emph{saka}
  \item \emph{Say it:} \emph{saka}, once more
  \item \emph{Say it:} say \emph{saka}
  \item \emph{Recall:} say the Japanese for foam, then the Japanese for dew; the rainy season, then say \emph{saka} again
\end{itemize}

\begin{tcolorbox}[breakable,colback=teal!4,colframe=teal!35!black,title={Before you move on}]
What does \ja{さか} mean? (\textquotedblleft{}A slope\textquotedblright{}.) Say it once more.
\end{tcolorbox}
